\chapter{Clase 3}

\section{Presión}

Cualquier fluido alcanza eventualmente un estado de equilibrio hidrostático
en ausencia de fuerzas impulsoras externas y ante la presencia de mecanismos
disipativos. En este estado, el movimiento cesa y los campos físicos dejan de
variar con el tiempo.

La presión es una fuerza de contacto que actúa sobre superficies reales o
imaginarias dentro de un cuerpo. Como un fluido en reposo no puede sostener
esfuerzos de cizalla, la fuerza de contacto actúa exclusivamente en dirección
normal a cada elemento de superficie.

\subsection{Definición y unidades}

La presión se define como la magnitud de la fuerza normal ejercida por unidad
de área:

$$
p=\frac{F_{\mathrm{n}}}{A}.
$$

En el Sistema Internacional, su unidad es el pascal,

$$
1\,\mathrm{Pa}=1\,\mathrm{N\,m^{-2}}.
$$

También se utilizan con frecuencia las siguientes unidades:

\begin{itemize}
    \item bar:
    $1\,\mathrm{bar}=10^5\,\mathrm{Pa}$;
    \item atmósfera estándar:
    $1\,\mathrm{atm}=101325\,\mathrm{Pa}$;
    \item hectopascal:
    $1\,\mathrm{hPa}=100\,\mathrm{Pa}$, equivalente a un milibar.
\end{itemize}

\subsection{El mar incompresible}

Considérese una columna vertical de un fluido en equilibrio mecánico, con
densidad constante $\rho_0$, área transversal $A$ y altura $h$. La presión en
la base es $p$, mientras que en la superficie superior es $p_0$. El balance de
fuerzas verticales exige que la fuerza ascendente en la base equilibre la
fuerza descendente en la superficie y el peso de la columna:

$$
pA=p_0A+Mg_0.
$$

Como la masa del fluido es

$$
M=\rho_0Ah,
$$

la presión en la base resulta

$$
p=p_0+\rho_0g_0h.
$$

Si el eje $z$ se orienta verticalmente hacia arriba y la superficie libre se
encuentra en $z=0$, los puntos bajo la superficie tienen $z<0$. En estas
coordenadas, el campo de presión es

$$
\boxed{p(z)=p_0-\rho_0g_0z}.
$$

Por tanto, la presión aumenta linealmente con la profundidad. En el mar, este
aumento es aproximadamente de un bar por cada diez metros de profundidad.

\subsubsection*{Ejemplo: fuerza sobre una compuerta}

Considérese una compuerta vertical de ancho $L$, con agua hasta una altura
$h$. La presión atmosférica $p_0$ actúa a ambos lados, por lo que la fuerza
neta depende del exceso de presión hidrostática:

$$
p(z)-p_0=-\rho_0g_0z.
$$

Un elemento horizontal de altura $dz$ tiene área $dA=L\,dz$. La fuerza total
sobre la compuerta es entonces

$$
\begin{aligned}
F
&=\int_{-h}^{0}\bigl(p(z)-p_0\bigr)L\,dz\\
&=-\rho_0g_0L\int_{-h}^{0}z\,dz\\
&=\frac{1}{2}\rho_0g_0Lh^2.
\end{aligned}
$$

Este resultado también puede interpretarse como el producto del área
$Lh$ por el exceso de presión promedio,
$\frac{1}{2}\rho_0g_0h$.

\subsubsection*{Ejemplo: barómetro de mercurio de Torricelli}

En un barómetro de mercurio, la columna alcanza una altura $h$ cuando su peso
equilibra la presión atmosférica exterior $p$. Si la presión $p_0$ en la parte
superior del tubo es prácticamente nula, se cumple

$$
p=\rho_0g_0h,
$$

y, por tanto,

$$
\boxed{h=\frac{p}{\rho_0g_0}}.
$$

Para una presión de una atmósfera y una densidad del mercurio de
$13.534\,\mathrm{g\,cm^{-3}}$, la altura de la columna es aproximadamente
$76.34\,\mathrm{cm}$.

\subsection{¿Atmósfera incompresible?}

Aunque el aire es compresible, el modelo de densidad constante proporciona
una aproximación razonable durante los primeros kilómetros de la atmósfera.
Si se supusiera una densidad constante $\rho_0$, la presión disminuiría
linealmente con la altura:

$$
p(z)=p_0-\rho_0g_0z.
$$

En este modelo, la presión se anularía a la altura característica

$$
h_0=\frac{p_0}{\rho_0g_0}.
$$

Para valores atmosféricos estándar cercanos a $25\,^{\circ}\mathrm{C}$, se
obtiene

$$
h_0\approx8.72\,\mathrm{km}.
$$

La presión real no se anula a esta altura porque la densidad del aire también
disminuye con $z$. Sin embargo, $h_0$ proporciona la escala vertical en la que
ocurren cambios atmosféricos importantes, mientras que la aproximación lineal
es útil aproximadamente durante los primeros $2\,\mathrm{km}$.

\subsection{Origen microscópico de la presión en gases}

En un gas, la presión surge del intercambio de momento producido por las
colisiones de las moléculas contra las paredes del recipiente. Considérese una
pared perpendicular al eje $z$ y una molécula de masa $m$ que se aproxima con
$v_z<0$. Si la colisión es elástica, su velocidad normal cambia de $v_z$ a
$-v_z$.

La magnitud del momento transferido a la pared por cada colisión es

$$
2m|v_z|.
$$

Si inicialmente se supone que todas las moléculas tienen la misma componente
$v_z<0$, la masa que alcanza un área $A$ durante un intervalo $dt$ es

$$
dM=\rho A(-v_z)\,dt.
$$

El momento transferido en dirección normal es

$$
dP_z=2v_z\,dM
=-2\rho v_z^2A\,dt,
$$

de modo que la fuerza ejercida sobre la pared es

$$
F_z=\frac{dP_z}{dt}
=-2\rho v_z^2A.
$$

El signo negativo indica que la fuerza apunta en dirección opuesta al eje
$z$. Para obtener la presión debe considerarse que solo la mitad de las
moléculas se mueve hacia la pared y que existe una distribución de
velocidades. La magnitud de la fuerza promedio por unidad de área es entonces

$$
p=\rho\langle v_z^2\rangle.
$$

En un gas isotrópico no existe una dirección privilegiada, por lo que

$$
\langle v_x^2\rangle
=\langle v_y^2\rangle
=\langle v_z^2\rangle
=\frac{1}{3}\langle v^2\rangle.
$$

Se obtiene así la relación microscópica

$$
\boxed{
p=\rho\langle v_z^2\rangle
=\frac{1}{3}\rho\langle v^2\rangle
}.
$$

La rapidez cuadrática media molecular se define como

$$
v_{\mathrm{mol}}=\sqrt{\langle v^2\rangle},
$$

por lo que

$$
\boxed{v_{\mathrm{mol}}=\sqrt{\frac{3p}{\rho}}}.
$$

Para el aire a una presión de una atmósfera y una temperatura cercana a
$18\,^{\circ}\mathrm{C}$, esta rapidez es aproximadamente
$500\,\mathrm{m\,s^{-1}}$, un valor del mismo orden de magnitud que la
velocidad del sonido.

\subsection{El campo de presión}

La presión se manifiesta como una fuerza de contacto que actúa sobre
superficies. Estas pueden ser interfaces físicas reales, donde las propiedades
del material cambian abruptamente, o superficies imaginarias que separan dos
partes de un mismo cuerpo.

\subsubsection{Definición matemática de la fuerza de presión}

Para analizar esta fuerza, se divide una superficie orientada en elementos
infinitesimales. El vector de superficie se define como

$$
d\vec{S}=\vec{n}\,dS,
$$

donde $dS$ es el área del elemento y $\vec{n}$ es el vector unitario normal.
Como un fluido en reposo no puede sostener fuerzas de cizalla, las fuerzas de
contacto actúan únicamente en la dirección normal.

La fuerza ejercida por el material situado en el lado positivo del elemento
sobre el material situado en el lado negativo se expresa mediante el campo de
presión $p(\vec{x})$ como

$$
d\vec{F}=-p(\vec{x})\,d\vec{S}.
$$

El signo negativo indica que una presión positiva produce una fuerza opuesta
a la normal exterior, es decir, una compresión. Una presión negativa
representaría una tensión sobre la superficie. La fuerza de presión total
sobre una superficie $S$ es

$$
\boxed{
\vec{F}=-\int_S p\,d\vec{S}
}.
$$

\subsubsection{Fuerza resultante sobre una partícula material}

Considérese una partícula material infinitesimal en forma de caja rectangular,
con lados $dx$, $dy$ y $dz$. Su volumen es

$$
dV=dx\,dy\,dz.
$$

La fuerza resultante surge de la pequeña diferencia de presión entre caras
opuestas. En la dirección $z$ se tiene

$$
dF_z
=\bigl[p(x,y,z)-p(x,y,z+dz)\bigr]dx\,dy.
$$

Al expandir la presión de la cara superior hasta primer orden,

$$
p(x,y,z+dz)
\approx p(x,y,z)+\frac{\partial p}{\partial z}dz,
$$

se obtiene

$$
dF_z
\approx-\frac{\partial p}{\partial z}dV.
$$

Al sumar las contribuciones de las tres direcciones espaciales, la fuerza de
presión total sobre la partícula es

$$
\boxed{
d\vec{F}=-\vec{\nabla}p\,dV
}.
$$

Esto significa que el efecto de todas las fuerzas de contacto es equivalente
a una fuerza de volumen cuya densidad es el gradiente negativo de la presión,
$-\vec{\nabla}p$.

\subsubsection{Teorema de Gauss}

La fuerza de presión total sobre un cuerpo de volumen $V$ y superficie cerrada
$S$ puede calcularse mediante una integral de superficie o una integral de
volumen. La equivalencia entre ambas expresiones conduce a

$$
\boxed{
\oint_S p\,d\vec{S}
=\int_V \vec{\nabla}p\,dV
}.
$$

Para un campo vectorial general $\vec{U}$, el teorema de Gauss se expresa como

$$
\boxed{
\oint_S \vec{U}\cdot d\vec{S}
=\int_V \vec{\nabla}\cdot\vec{U}\,dV
}.
$$

La integral de superficie representa el flujo de $\vec{U}$ a través de la
frontera, mientras que la integral de volumen suma su divergencia en el
interior del cuerpo.

\subsubsection{Ley de Pascal: isotropía de la presión}

La presión en un fluido en reposo no depende de la orientación del elemento de
superficie. Esta propiedad se conoce como ley de Pascal y establece que la
presión es un campo escalar.

Para demostrarlo, considérese una partícula material con forma de tetraedro
recto orientado según los ejes coordenados. Sean $dS$ el área de la cara
oblicua y $dS_x$, $dS_y$ y $dS_z$ las áreas de sus proyecciones sobre los
planos coordenados. Si las presiones en las caras fueran $p$, $p_x$, $p_y$ y
$p_z$, el equilibrio mecánico exigiría

$$
d\vec{F}_S+d\vec{F}_V=\vec{0},
$$

donde $d\vec{F}_S$ es la fuerza total de contacto y

$$
d\vec{F}_V=\vec{f}\,dV
$$

es una fuerza de volumen, como la gravedad.

Si todas las dimensiones del tetraedro se reducen por un factor
$\lambda<1$, las áreas disminuyen como $\lambda^2$, mientras que el volumen y
las fuerzas de volumen disminuyen como $\lambda^3$. En el límite
$\lambda\to0$, la fuerza de volumen se vuelve despreciable frente a las
fuerzas de superficie. Por consiguiente,

$$
d\vec{F}_S=\vec{0}.
$$

El equilibrio de las componentes de esta fuerza, junto con las relaciones
geométricas entre la cara oblicua y sus proyecciones, exige que

$$
\boxed{p_x=p_y=p_z=p}.
$$

Por tanto, la presión en un punto de un fluido en reposo es la misma en todas
las direcciones.
